% ============================================================
% sec/03baseline_literatura.tex
%
% DEBE SALIR DEL EJE C. Un NÚMERO OBJETIVO con su fuente.
%
% DEBE CONTENER:
%   - El mejor resultado publicado con la métrica elegida,
%     o un rango razonable si la comparación directa no es posible
%   - Declaración EXPLÍCITA de qué se busca igualar o superar
%   - Las diferencias de protocolo que impiden una comparación justa
%
% Si el dataset carece de literatura previa, el requisito se
% sustituye por datasets análogos del mismo dominio, justificando
% por qué son o no comparables.
%
% Este número se arrastra a la Etapa 2 como tercera referencia,
% junto a los baselines internos.
% ============================================================
\section{Baseline de Literatura}
\label{sec:baseline}

\subsection{Origen del baseline}
\label{subsec:origenbaseline}
\pc{De qué trabajo sale el número, sobre qué datos, con qué métrica y
bajo qué protocolo. Si es un dataset análogo y no el propio, decirlo
aquí sin rodeos.}

\subsection{Número objetivo}
\label{subsec:objetivobaseline}
\begin{quote}
\textbf{Objetivo declarado:} \pc{Valor + métrica + fuente. Si la
comparación directa no es posible, declarar un rango razonable y
explicar por qué es un rango y no un número.}
\end{quote}

\subsection{Comparabilidad del protocolo}
\label{subsec:comparabilidad}
\pc{Un párrafo declarando el estatus de la comparación: equivalente,
parcialmente comparable, o solo indicativa del orden de magnitud.}

\noindent\textbf{Diferencias que impiden una comparación justa.}
\begin{itemize}
  \item \pc{Partición distinta}
  \item \pc{Métrica o variante de métrica distinta}
  \item \pc{Subconjunto de datos distinto}
  \item \pc{Población, dominio o reglas distintas}
\end{itemize}

\noindent\textbf{Qué significa entonces el número.} \pc{Declarar si es
una cifra a batir o solo una referencia del orden de magnitud de lo que
se considera un buen resultado en esta tarea. Presentarlo como
benchmark equivalente cuando no lo es sería una atribución incorrecta.}